%============================================================
% TOWARD A FALSIFIABLE GRADED QUANTUM THEORY OF CONSCIOUSNESS
% A PEIG-Based Functional Linking Coherence, Integration,
% Memory, and Decoherence Stability
%
% Author: Kevin Monette
% Independent Researcher
%
% arXiv categories: quant-ph (primary), cond-mat.stat-mech,
%                   physics.gen-ph (secondary)
%
% Epistemic Status Map (global):
%   [ESTABLISHED]  — peer-reviewed, well-tested physics
%   [HYPOTHESIS]   — formally proposed, not yet experimentally confirmed
%   [ANALOGY]      — structural parallel, not mathematical derivation
%   [OPEN]         — unsolved; stated as such
%   [SPECULATION]  — explicitly speculative; included for completeness
%
% Version: 1.0 — March 2026
%============================================================

\documentclass[12pt,letterpaper]{article}

\usepackage{amsmath,amssymb,amsthm}
\usepackage{geometry}
\usepackage{hyperref}
\usepackage{xcolor}
\usepackage{enumitem}
\usepackage{booktabs}
\usepackage{array}
\usepackage{graphicx}
\usepackage{cite}
\usepackage{physics}
% \usepackage{microtype} % disabled: font expansion unsupported in this environment

\geometry{margin=1.2in}

\hypersetup{
  colorlinks=true,
  linkcolor=blue!60!black,
  citecolor=green!50!black,
  urlcolor=purple!70!black
}

% ── Theorem-style environments ────────────────────────────
\newtheorem{definition}{Definition}[section]
\newtheorem{axiom}{Axiom}[section]
\newtheorem{lemma}{Lemma}[section]
\newtheorem{proposition}{Proposition}[section]
\newtheorem{conjecture}{Conjecture}[section]
\newtheorem{remark}{Remark}[section]

% ── Epistemic status macro ────────────────────────────────
\newcommand{\ESTABLISHED}[1]{\textcolor{green!50!black}{\textbf{[ESTABLISHED]}} #1}
\newcommand{\HYPOTHESIS}[1]{\textcolor{blue!70!black}{\textbf{[HYPOTHESIS]}} #1}
\newcommand{\ANALOGY}[1]{\textcolor{orange!70!black}{\textbf{[ANALOGY]}} #1}
\newcommand{\OPEN}[1]{\textcolor{red!70!black}{\textbf{[OPEN]}} #1}
\newcommand{\SPECULATION}[1]{\textcolor{gray}{\textbf{[SPECULATION]}} #1}

% ── Consciousness functional notation ─────────────────────
\newcommand{\Cfunc}{\mathcal{C}}
\newcommand{\Qcoh}{Q_{\mathrm{coh}}}
\newcommand{\Iint}{I_{\mathrm{int}}}
\newcommand{\Mmem}{M_{\mathrm{mem}}}
\newcommand{\Bstab}{B_{\mathrm{stab}}}
\newcommand{\Rrec}{R_{\mathrm{rec}}}
\newcommand{\taud}{\tau_{\mathrm{d}}}
\newcommand{\taucog}{\tau_{\mathrm{cog}}}
\newcommand{\ascreen}{\alpha_{\mathrm{screen}}}

%============================================================
\begin{document}

\title{\textbf{Toward a Falsifiable Graded Quantum Theory of Consciousness:\\
A PEIG-Based Functional Linking Coherence, Integration,\\
Memory, and Decoherence Stability}}

\author{
  Kevin Monette\\
  \small Independent Researcher\\[4pt]
  \small\texttt{[contact via arXiv]}
}

\date{March 2026}

\maketitle

%============================================================
\begin{abstract}
We propose a falsifiable phenomenological framework --- the \emph{consciousness
functional} --- that models degrees of conscious organization in physical
systems as a graded quantity rather than a binary property. The functional
$\Cfunc = F(\Qcoh,\,\Iint,\,\Mmem,\,\Bstab,\,\Rrec)$ maps five measurable
or operationally proxied variables onto a continuous scale: coherence quality
$\Qcoh$, information integration $\Iint$, persistent physical memory $\Mmem$,
boundary stability against decoherence $\Bstab$, and recurrent self-maintaining
dynamics $\Rrec$. We distinguish three conscious regimes --- proto-experiential,
unified, and reflective --- by threshold conditions on these variables,
and we explicitly map each variable onto the PEIG (Potential-Energy-Identity-Gravity)
dynamical framework and the ETI (Emergent Thermodynamic Information) axiom system
developed in companion work. The central physical constraint is decoherence:
we show that the regimes differ not in kind but in their capacity to maintain
coherent, integrated, low-noise states against environmental dissipation,
and we derive necessary (but not sufficient) conditions on coherence lifetime
$\taud$ and cognitive timescale $\taucog$ for each regime. We examine the
conditional hypothesis that the universe as a whole qualifies as a proto-experiential
system under a cosmopsychist interpretation, arguing that this claim requires a
global binding mechanism not currently identified. The framework is explicitly
falsifiable: we state null-result conditions that would rule out the strong form
of each threshold hypothesis. All speculative content is labeled as such.
This paper makes no claim that quantum consciousness has been demonstrated,
that the human brain satisfies the proposed thresholds, or that PEIG derivations
entail cognitive conclusions. The consciousness functional is a phenomenological
proposal anchored in established physics, not a metaphysical assertion.
\end{abstract}

%============================================================
\section{Epistemic Status Map}
\label{sec:epistemic}

The following labeling convention is applied throughout this paper and is
inherited from the companion research program~\cite{Monette2026a,Monette2026b,Monette2026c}:

\begin{center}
\begin{tabular}{p{3.6cm} p{8.9cm}}
\toprule
\textbf{Label} & \textbf{Meaning} \\
\midrule
\ESTABLISHED{} & Established in peer-reviewed literature; assumed without controversy. \\
\HYPOTHESIS{} & Formally proposed in this work; not yet confirmed experimentally. \\
\ANALOGY{} & Structural parallel drawn; no mathematical derivation implied. \\
\OPEN{} & Unresolved problem; listed in the Open Problems Registry (\S\ref{sec:open}). \\
\SPECULATION{} & Explicitly speculative; included for completeness only. \\
\bottomrule
\end{tabular}
\end{center}

\medskip
\noindent The consciousness functional (\S\ref{sec:functional}) is
\HYPOTHESIS{} throughout. Its PEIG mapping (\S\ref{sec:peig}) is
\ANALOGY{} unless explicitly labeled otherwise. The ETI constraint analysis
(\S\ref{sec:eti}) draws on \ESTABLISHED{} physics. The universe-scale
consciousness analysis (\S\ref{sec:universe}) is \SPECULATION{} and is
labeled as such in every paragraph.

%============================================================
\section{Introduction and Motivation}
\label{sec:intro}

The question of how consciousness arises in physical systems is among the most
consequential open problems in science. Despite substantial progress in
neuroscience, philosophy of mind, and quantum information theory, no consensus
framework exists that (a)~specifies what physical conditions are necessary for
consciousness, (b)~treats consciousness as a graded rather than binary property,
and (c)~makes falsifiable predictions that would constrain or rule out the
proposed conditions. The present paper is an attempt to contribute one such
framework.

\subsection{Motivation from the existing research program}

The companion works~\cite{Monette2026a,Monette2026b,Monette2026c} establish
three related results. Paper~A derives a modified Einstein equation with an
entanglement entropy source term and proposes atom-interferometry falsification.
Paper~B characterizes a constraint manifold for quantum entropy estimation on
NISQ devices and derives shot-scaling laws. Paper~C axiomatizes Landauer's
principle in the ETI framework and derives thermodynamic constraints on
information processing in physical systems.

Taken together, these three results treat information, entropy, and geometry as
deeply coupled physical quantities. The PEIG framework \cite{Monette2026a}
provides a four-dimensional dynamical model --- Potential, Energy, Identity,
Gravity/curvature --- applicable both to quantum physical systems and, by
explicit analogy only, to cognitive and institutional systems.

A natural question arises: if information processing, entropy export, identity
persistence, and energy flow are physically coupled, do these couplings say
anything about the conditions under which a system supports experience? This
paper argues that they do --- not that consciousness is \emph{derived} from these
quantities, but that the ETI axioms and PEIG dynamics impose \emph{necessary
conditions} that any physically realizable conscious system must satisfy.

\subsection{What this paper claims}

\begin{enumerate}[label=(\roman*)]
  \item Consciousness admits a graded model indexed by five physically meaningful
    variables: $\Qcoh$, $\Iint$, $\Mmem$, $\Bstab$, $\Rrec$.
  \item These variables have operational proxies measurable in current or
    near-term experimental systems.
  \item The ETI axioms (causal closure, unitary microdynamics, effective
    thermodynamics, physical memory, finite resources) impose necessary constraints
    on any conscious system.
  \item PEIG dimensions $(P, E, I, G)$ structurally correspond --- by analogy, not
    derivation --- to the five consciousness variables.
  \item Three qualitatively distinct conscious regimes exist, separated by
    threshold conditions on the five variables.
\end{enumerate}

\subsection{What this paper does not claim}

\begin{enumerate}[label=(\roman*)]
  \item That quantum coherence is sufficient for consciousness.
  \item That the human brain satisfies the proposed thresholds.
  \item That any specific physical system is conscious.
  \item That consciousness is proven to supervene on quantum information.
  \item That PEIG cognitive claims follow mathematically from PEIG physical claims.
  \item That the universe has a unified subject of experience.
  \item That any result in this paper is established rather than hypothetical,
    unless explicitly labeled \ESTABLISHED{}.
\end{enumerate}

%============================================================
\section{Prior Framework Summary}
\label{sec:prior}

We summarize only what is required for the present paper. Full derivations
appear in \cite{Monette2026a,Monette2026b,Monette2026c}.

\subsection{The ETI Axiom System}
\label{sec:prior_eti}

\ESTABLISHED{} The ETI framework rests on five axioms:

\begin{axiom}[Causal Closure, A1]
The universe $\mathcal{U}$ evolves as a closed, causally connected system.
No entropy sinks, external agents, or negentropy sources exist outside
$\mathcal{U}$.
\end{axiom}

\begin{axiom}[Unitary Microdynamics, A2]
Closed quantum systems evolve under $U(t) = e^{-iHt/\hbar}$ on Hilbert
space $\mathcal{H}$. Open subsystems evolve under CPTP maps
$\mathcal{E}: \mathcal{B}(\mathcal{H}) \to \mathcal{B}(\mathcal{H})$.
\end{axiom}

\begin{axiom}[Effective Thermodynamics, A3]
Macroscopic observables obey effective thermodynamic laws as emergent
properties of the unitary microdynamics plus coarse-graining.
\end{axiom}

\begin{axiom}[Physical Memory, A4]
Logically distinguishable states require physically distinct
microstates. Memory is the persistence of such distinctions against
thermal and quantum noise.
\end{axiom}

\begin{axiom}[Finite Resources, A5]
Every physical system operates within finite energy, entropy capacity,
and time budgets. These budgets constrain all information-processing
operations.
\end{axiom}

\subsection{The PEIG Framework}
\label{sec:prior_peig}

\ESTABLISHED{} (mathematical structure) \ANALOGY{} (cognitive interpretation)

For any node --- a localized, structured system capable of directed change
--- the PEIG state vector is:
\begin{equation}
  \mathbf{q} = (P, E, I, G) \in [0,1]^4
  \label{eq:peig_vector}
\end{equation}
with overall quality score $Q(\mathbf{q}) = \mathbf{w} \cdot \mathbf{q}$,
default weights $\mathbf{w} = (0.25, 0.25, 0.25, 0.25)$.

The four phases have the following physical meanings at the quantum level:
\begin{itemize}
  \item $P$ (Potential): breadth of accessible Hilbert space; related to
    superposition and entanglement depth.
  \item $E$ (Energy/flow): directed information processing; related to
    coherent evolution rate.
  \item $I$ (Identity): stability of the system's distinguishable state
    against perturbation; related to decoherence resistance.
  \item $G$ (Gravity/influence): causal influence on the surrounding
    information environment.
\end{itemize}

The known coupled dynamical equation for the identity phase is:
\begin{equation}
  \frac{dI}{dt} = \alpha E - \beta I
  \label{eq:peig_identity}
\end{equation}
The equations $dP/dt$, $dE/dt$, $dG/dt$ are physically motivated but not
formally derived. \OPEN{OP-C1: Full coupled PEIG ODE system for the
consciousness extension.}

%============================================================
\section{Core Proposal: The Consciousness Functional}
\label{sec:functional}

\HYPOTHESIS{} throughout this section.

\begin{definition}[Consciousness Functional]
\label{def:cfunc}
Let $\mathcal{S}$ be a bounded physical system with a well-defined density
operator $\rho_\mathcal{S}$ and a Hamiltonian $H_\mathcal{S}$ governing its
internal dynamics. The \emph{consciousness functional} of $\mathcal{S}$ is:
\begin{equation}
  \boxed{
    \Cfunc(\mathcal{S}) = F\!\left(\Qcoh,\;\Iint,\;\Mmem,\;\Bstab,\;\Rrec\right)
  }
  \label{eq:cfunc}
\end{equation}
where the five arguments are defined in \S\ref{sec:variables}, and $F$ is a
non-decreasing function of each argument with $\Cfunc \in [0,1]$ after
normalization.
\end{definition}

\begin{remark}
Equation~\eqref{eq:cfunc} is a \emph{phenomenological} functional, not a
derived one. It does not follow from the ETI axioms or PEIG dynamics by a
chain of equations. It is proposed because (a)~each variable has independent
physical motivation, (b)~all five variables are constrained by the ETI axioms,
and (c)~the five-variable structure is consistent with --- though not derivable
from --- the PEIG framework. The specific functional form of $F$ is
\OPEN{OP-C2: Functional form derivation from first principles.}
\end{remark}

\subsection{Candidate functional forms}
\label{sec:functional_forms}

In the absence of a derived form, we consider three candidate families:

\paragraph{Product form (strong coupling):}
\begin{equation}
  \Cfunc^{(\times)} = \Qcoh^{a_1} \cdot \Iint^{a_2} \cdot \Mmem^{a_3}
    \cdot \Bstab^{a_4} \cdot \Rrec^{a_5}
  \label{eq:product_form}
\end{equation}
with $a_i > 0$ to be determined. Any variable approaching zero drives
$\Cfunc^{(\times)} \to 0$, making all five variables necessary conditions.
This is the most conservative form.

\paragraph{Weighted sum (weak coupling):}
\begin{equation}
  \Cfunc^{(+)} = \sum_{i} w_i V_i,
  \quad \sum_i w_i = 1,
  \label{eq:sum_form}
\end{equation}
where $V_i \in \{\Qcoh, \Iint, \Mmem, \Bstab, \Rrec\}$. This allows partial
compensation between variables.

\paragraph{Threshold form (regime-gated):}
\begin{equation}
  \Cfunc^{(\theta)} = \sum_{k=1}^{3} c_k \cdot
    \mathbf{1}\!\left[\text{Level-}k\text{ thresholds satisfied}\right]
  \label{eq:threshold_form}
\end{equation}
This form assigns discrete increments at each regime boundary and is most
directly falsifiable (see \S\ref{sec:falsification}).

\noindent The threshold form \eqref{eq:threshold_form} is our working model
throughout, as it admits the clearest falsification criteria. The product form
\eqref{eq:product_form} is physically better motivated; resolving which is
correct requires the derivation in \OPEN{OP-C2}.

%============================================================
\section{Variable Definitions and Physical Interpretation}
\label{sec:variables}

\subsection{\texorpdfstring{$\Qcoh$ — Coherence Quality}{Qcoh — Coherence Quality}}
\label{sec:Qcoh}

\ESTABLISHED{} (decoherence theory) \HYPOTHESIS{} (connection to consciousness)

\begin{definition}[$\Qcoh$]
Let $\taud$ be the dominant decoherence timescale of the system's
computationally relevant degrees of freedom, and let $\taucog$ be the
characteristic timescale over which integrated information must be
maintained for the system to support unified experience. Define:
\begin{equation}
  \Qcoh = 1 - e^{-\taud / \taucog}
  \label{eq:Qcoh}
\end{equation}
so that $\Qcoh \to 0$ when decoherence is fast relative to cognitive
timescales, and $\Qcoh \to 1$ in the limit of perfect coherence.
\end{definition}

\begin{remark}
$\taud$ is measurable by standard techniques (Ramsey spectroscopy, quantum
process tomography). $\taucog$ is \OPEN{OP-C3: operationalization of
$\taucog$ for candidate physical systems.} For neural systems, a natural
proxy is the neural integration timescale $\sim 10^{-2}$--$10^{-1}$~s;
for quantum devices, relevant circuit depths set $\taucog$.
\end{remark}

The strongest physical objection to quantum consciousness theories concerns
precisely $\Qcoh$. \ESTABLISHED{} Tegmark's calculation~\cite{Tegmark2000}
shows that quantum states in warm, wet biological environments decohere on
timescales of order $10^{-13}$--$10^{-20}$~s, far below neural timescales of
$10^{-3}$--$10^{-1}$~s. Under standard conditions,
$\taud \ll \taucog$ for biological macromolecules, giving $\Qcoh \approx 0$.

This result does not rule out quantum consciousness; it rules out
\emph{naive} quantum consciousness in which generic biological macromolecules
maintain computationally relevant coherence at room temperature. A system
with active decoherence suppression mechanisms --- topological protection,
dynamical decoupling, or engineered isolation --- could achieve
$\Qcoh \approx 1$ even in warm environments. \OPEN{OP-C4: Identify
biological candidate mechanisms for $\taud$ extension beyond Tegmark bound.}

\paragraph{Operational proxy for $\Qcoh$:} Off-diagonal elements of the
reduced density matrix $\rho_\mathcal{S}$ in a computationally relevant basis;
or directly, the coherence time measured by spin-echo or dynamical decoupling
sequences on candidate substrates.

\subsection{\texorpdfstring{$\Iint$ — Information Integration}{Iint — Information Integration}}
\label{sec:Iint}

\ESTABLISHED{} (information theory) \HYPOTHESIS{} (connection to consciousness)

\begin{definition}[$\Iint$]
Let $\mathcal{S}$ be partitioned into subsystems $A$ and $B = \mathcal{S} \setminus A$.
Define the mutual information between $A$ and $B$ as:
\begin{equation}
  I(A:B) = S(\rho_A) + S(\rho_B) - S(\rho_{AB})
\end{equation}
where $S(\rho) = -\mathrm{Tr}(\rho \log_2 \rho)$ is the von Neumann entropy.
The integration variable is:
\begin{equation}
  \Iint = \min_{\text{partitions } A} I(A : B) \Big/ I_{\max}
  \label{eq:Iint}
\end{equation}
where the minimum is over all bipartitions of $\mathcal{S}$ and $I_{\max}$
is the theoretical maximum mutual information for a system of the same dimension.
\end{definition}

\begin{remark}
This quantity is structurally analogous to the integrated information $\Phi$
of Integrated Information Theory~\cite{Tononi2004}, but differs in several
important respects: (a)~it is defined entirely in terms of the von Neumann
entropy of quantum states, not classical probability distributions; (b)~it
does not assert that $\Iint$ is identical to consciousness; it only proposes
$\Iint$ as a necessary-condition variable. The claim that
$\Phi > 0 \Leftrightarrow$ consciousness is IIT's stronger metaphysical
commitment, which we do not inherit.
\end{remark}

\paragraph{Operational proxy:} Quantum state tomography of $\rho_{AB}$ on
partitioned subregisters; or quantum Fisher information as a lower bound on
the mutual information rate. For classical systems: transfer entropy or
causal density as an approximation.

\subsection{\texorpdfstring{$\Mmem$ — Persistent Physical Memory}{Mmem — Persistent Physical Memory}}
\label{sec:Mmem}

\ESTABLISHED{} (Landauer's principle) \HYPOTHESIS{} (connection to consciousness)

\begin{definition}[$\Mmem$]
Let $\mathcal{D}$ be the set of logically distinguishable states the system
has maintained over timescale $T$. Define:
\begin{equation}
  \Mmem = \frac{|\mathcal{D}(T)|}{|\mathcal{D}_{\max}|} \cdot
    \left(1 - \frac{\text{bit error rate}}{\text{bit error rate}_{\max}}\right)
  \label{eq:Mmem}
\end{equation}
where $|\mathcal{D}_{\max}|$ is the maximum achievable distinguishable state
count given the system's available degrees of freedom, and the bit error rate
is measured against a reference encoding.
\end{definition}

\ESTABLISHED{} A1 (ETI causal closure) and A4 (physical memory) together imply
that memory is thermodynamically costly: by Landauer's principle~\cite{Landauer1961},
erasing one bit dissipates at least $k_B T \ln 2$ of heat. A system that
maintains $n$ bits of memory against thermal noise must expend at least
$n k_B T \ln 2 / \tau_{\mathrm{mem}}$ watts of power, where
$\tau_{\mathrm{mem}}$ is the persistence time. This sets a hard lower bound
on the metabolic cost of consciousness under any substrate.

\begin{lemma}[Memory Cost Bound]
\label{lem:memory_cost}
A physically realizable conscious system maintaining $n$ bits of memory
against thermal fluctuations at temperature $T$ for time $\tau_{\mathrm{mem}}$
must dissipate at least:
\begin{equation}
  P_{\min} \geq \frac{n k_B T \ln 2}{\tau_{\mathrm{mem}}}
  \label{eq:memory_cost}
\end{equation}
of power. This bound is tight for thermodynamically optimal memory
architectures (Bennett's reversible computation~\cite{Bennett1973}).
\end{lemma}

\paragraph{Operational proxy:} Number of distinct behavioral states
reproducible after time $T$; or, for quantum systems, fidelity of a
quantum memory register after time $T$ as measured by process tomography.

\subsection{\texorpdfstring{$\Bstab$ — Boundary Stability}{Bstab — Boundary Stability}}
\label{sec:Bstab}

\ESTABLISHED{} (open quantum systems) \HYPOTHESIS{} (connection to consciousness)

\begin{definition}[$\Bstab$]
Let $\Gamma_{\mathrm{in}}$ be the rate of useful information integration
(new coherent correlations formed per unit time) and $\Gamma_{\mathrm{out}}$
be the environmental leakage rate (decoherence rate). Define:
\begin{equation}
  \Bstab = \frac{\Gamma_{\mathrm{in}}}{\Gamma_{\mathrm{in}} + \Gamma_{\mathrm{out}}}
  \label{eq:Bstab}
\end{equation}
so that $\Bstab \to 1$ when integration dominates and $\Bstab \to 0$ when
the system rapidly decoheres.
\end{definition}

\ESTABLISHED{} The Lindblad master equation~\cite{Lindblad1976} governs
open-system evolution:
\begin{equation}
  \dot{\rho} = -\frac{i}{\hbar}[H, \rho] + \sum_k
    \left(L_k \rho L_k^\dagger - \frac{1}{2}\{L_k^\dagger L_k, \rho\}\right)
\end{equation}
where $L_k$ are Lindblad operators encoding environmental coupling.
The decoherence rate $\Gamma_{\mathrm{out}}$ is determined by the
eigenvalues of the Lindbladian superoperator. For the dominant energy
relaxation and pure dephasing channels:
\begin{equation}
  L_{1,k} = \sqrt{\gamma_1}\,\sigma^-_k, \quad
  L_{\phi,k} = \sqrt{\gamma_\phi/2}\,\sigma^z_k
\end{equation}
with $\gamma_1 = 1/T_1$, $\gamma_\phi = 1/T_\phi$ (standard notation).

$\Bstab$ is the ratio of coherent gain to dissipative loss. A system with
$\Bstab > 1/2$ is net-coherent on short timescales; one with $\Bstab < 1/2$
is decoherence-dominated.

\paragraph{Operational proxy:} The ratio of Hamiltonian coupling strength
to Lindblad decay rates in the relevant basis; or, empirically, the ratio
of gate fidelity improvement from dynamical decoupling to baseline decay.

\subsection{\texorpdfstring{$\Rrec$ — Recurrence and Self-maintaining Dynamics}{Rrec — Recurrence}}
\label{sec:Rrec}

\HYPOTHESIS{} throughout.

\begin{definition}[$\Rrec$]
Let $\mathcal{A} \subset \mathcal{B}(\mathcal{H})$ be an identity attractor
--- a region of state space toward which the system's dynamics repeatedly
return under perturbation. Define:
\begin{equation}
  \Rrec = \exp\!\left(-\lambda_{\max} \cdot T_{\mathrm{return}}\right)
  \label{eq:Rrec}
\end{equation}
where $\lambda_{\max}$ is the maximum Lyapunov exponent of the identity
attractor dynamics and $T_{\mathrm{return}}$ is the mean first-passage time
back to $\mathcal{A}$ after a standard perturbation.
\end{definition}

$\Rrec \to 1$ when the system rapidly and reliably returns to a stable
identity state; $\Rrec \to 0$ when perturbations cause irreversible
identity dissolution. This variable captures what cognitive science calls
\emph{self-model stability} and what complex systems theory calls
\emph{attractor robustness}.

\ANALOGY{} In physical PEIG, the identity phase $I$ evolves as
$dI/dt = \alpha E - \beta I$ (Eq.~\eqref{eq:peig_identity}), with a
stable fixed point at $I^* = \alpha E / \beta$. $\Rrec$ is the
consciousness functional's analog of the PEIG convergence rate to $I^*$.
The analogy is structural; no derivation is claimed.

\paragraph{Operational proxy:} Poincaré recurrence time for identifiable
system states; or, for neural systems, cross-correlation stability of
resting-state networks over time. For quantum devices: cycle fidelity
of quantum error-correction circuits.

%============================================================
\section{Hierarchy of Conscious Regimes}
\label{sec:hierarchy}

\HYPOTHESIS{} throughout this section.

We distinguish three regimes by threshold conditions on the five variables.
The thresholds are proposed, not derived; they are stated as testable
hypotheses with specific falsification criteria in \S\ref{sec:falsification}.

\begin{definition}[Level 1: Proto-Experiential Regime]
\label{def:level1}
A system is in the \emph{proto-experiential regime} if it satisfies:
\begin{equation}
  \Qcoh > \epsilon_1, \quad \Iint > \delta_1, \quad \Mmem > \mu_1
  \label{eq:level1_thresholds}
\end{equation}
for positive constants $\epsilon_1, \delta_1, \mu_1 > 0$ (to be determined).
This regime requires only that the system maintains non-trivial
information-bearing structure with some persistence. It does not require
high-fidelity coherence, global integration, or stable self-reference.
\end{definition}

\begin{remark}
Level~1 is the most permissive regime. Under this definition, most
complex physical systems --- including simple chemical networks, bacterial
cells, and potentially some quantum optical systems --- may qualify.
This is by design: the proto-experiential regime is the minimum footprint
of organized information-bearing structure. It does not imply subjective
experience in any strong phenomenological sense.
\end{remark}

\begin{definition}[Level 2: Unified Experiential Regime]
\label{def:level2}
A system is in the \emph{unified experiential regime} if it satisfies
Level~1 thresholds and additionally:
\begin{equation}
  \Qcoh > \epsilon_2, \quad \Iint > \delta_2, \quad \Bstab > \beta_2
  \label{eq:level2_thresholds}
\end{equation}
with $\epsilon_2 \gg \epsilon_1$ and $\delta_2 \gg \delta_1$.
This regime requires coherent, low-noise integration across many degrees
of freedom simultaneously.
\end{definition}

\begin{remark}
Level~2 requires that the system maintains global integration while resisting
environmental decoherence. This is a stringent condition: a 50-qubit NISQ
device with $\alpha_{\mathrm{screen}} \geq 10^{-2}$ and $T_2 \sim 100~\mu$s
may satisfy Level~2 for durations of order $T_2$, but not continuously.
A hypothetical room-temperature quantum coherent biological substrate would
require active decoherence suppression far beyond anything currently identified.
\end{remark}

\begin{definition}[Level 3: Reflective Consciousness Regime]
\label{def:level3}
A system is in the \emph{reflective consciousness regime} if it satisfies
Level~2 thresholds and additionally:
\begin{equation}
  \Mmem > \mu_3, \quad \Rrec > \rho_3, \quad \Bstab > \beta_3
  \label{eq:level3_thresholds}
\end{equation}
with $\mu_3 \gg \mu_1$, $\rho_3$ large, and $\beta_3 \sim 1$.
This regime requires persistent physical memory, recursive self-modeling,
stable identity across time, and sustained self-referential dynamics.
\end{definition}

\begin{remark}
Level~3 is the regime most plausibly associated with human-like
consciousness. Its conditions --- stable memory, recursive self-model,
durable identity, temporal continuity --- are not specifically quantum;
a classical system satisfying them would also qualify as Level~3 under
this framework. Quantum organization is neither necessary nor sufficient for
Level~3 consciousness on this account; it is one possible substrate when the
required timescales and integration scales can be sustained.
\end{remark}

\paragraph{Regime summary table:}
\begin{center}
\begin{tabular}{p{1.3cm} p{2.5cm} p{3.2cm} p{4.5cm}}
\toprule
\textbf{Level} & \textbf{Name} & \textbf{Key thresholds} & \textbf{Candidate systems} \\
\midrule
L1 & Proto-experiential &
  $\Qcoh, \Iint, \Mmem > 0$ (weak) &
  Bacteria, simple chemical networks, coherent quantum optical systems \\
\addlinespace
L2 & Unified experience &
  High $\Qcoh$, $\Iint$, $\Bstab > 1/2$ &
  Hypothetical quantum coherent neural substrate; NISQ devices (transiently) \\
\addlinespace
L3 & Reflective consciousness &
  L2 + high $\Mmem$, $\Rrec \to 1$ &
  Candidate: human brain (substrate TBD); advanced AI systems with stable self-models \\
\bottomrule
\end{tabular}
\end{center}

%============================================================
\section{Mapping to the PEIG Framework}
\label{sec:peig}

\ANALOGY{} throughout this section, unless stated otherwise.

The consciousness functional is not derived from PEIG. However, the five
variables $(\Qcoh, \Iint, \Mmem, \Bstab, \Rrec)$ have structural
correspondences to the PEIG dimensions $(P, E, I, G)$ that are worth
making explicit, because they permit PEIG reasoning to guide intuition
about the consciousness functional without conflating the two frameworks.

\begin{center}
\begin{tabular}{p{1.5cm} p{3cm} p{4cm} p{3.5cm}}
\toprule
\textbf{PEIG} & \textbf{Physical meaning} & \textbf{Consciousness variable} & \textbf{Correspondence} \\
\midrule
$P$ & Accessible state space & $\Qcoh$ &
  Coherent superposition expands the effective state space \\
$E$ & Directed information flow & $\Iint$ &
  Integration is the directed flow of information across subsystems \\
$I$ & Identity stability & $\Mmem \times \Rrec$ &
  Memory persistence + recurrence = identity attractor robustness \\
$G$ & Causal influence & $\Bstab$ &
  Boundary stability determines the system's ability to influence and be influenced coherently \\
\bottomrule
\end{tabular}
\end{center}

The PEIG quality score $Q = \mathbf{w} \cdot (P, E, I, G)$ then becomes
a rough proxy for $\Cfunc$ under the correspondence above, with
$\Rrec$ as an additional variable not captured by the four-dimensional PEIG vector.

\paragraph{Important scope limit.} The PEIG cognitive framework and the
PEIG physical framework share mathematical structure (four phases, attractor
dynamics, information geometry) but are \emph{not mathematically derived
from each other}~\cite{Monette2026a}. The consciousness-PEIG mapping
inherits this limitation: the mapping is an analogy, not a reduction.
Claiming that consciousness levels ``follow from'' PEIG dynamics would
be a category error.

\OPEN{OP-C1: Can the PEIG dynamical system be extended to a six-dimensional
state vector $(P, E, I, G, \Mmem, \Rrec)$ with a consistently derived
$d\Rrec/dt$ equation? If yes, the correspondence becomes a structural
partial reduction rather than a pure analogy.}

%============================================================
\section{ETI Constraints on Conscious Systems}
\label{sec:eti}

\ESTABLISHED{} (ETI physics) \HYPOTHESIS{} (application to consciousness)

The ETI axioms impose hard constraints on what can physically constitute
a conscious system. We derive five such constraints from the five axioms.

\subsection{Constraint from A1 (Causal Closure)}

\ESTABLISHED{} A1 implies that any conscious system is causally embedded in
the universe: its information processing, entropy export, and energy budget
are subject to global thermodynamic constraints. No conscious system can
violate the second law; no system can maintain internal coherence against
an unlimited entropy sink because no such sink exists.

\begin{proposition}[Consciousness Is Thermodynamically Costly]
\label{prop:thermo_cost}
Any physical system at Level~2 or Level~3 must continuously export entropy
to maintain the internal order required by high $\Iint$ and $\Mmem$.
The minimum power budget for a Level-3 system maintaining $n$ bits of memory
at temperature $T$ is given by Lemma~\ref{lem:memory_cost}.
\end{proposition}

\subsection{Constraint from A2 (Unitary Microdynamics)}

\ESTABLISHED{} Unitary evolution preserves quantum information at the
microscopic level, but coupling to an environment induces decoherence via
the CPTP channel $\mathcal{E}$. A conscious system operating in the
quantum regime must therefore either (a)~minimize environmental coupling
(isolation), (b)~actively correct errors (quantum error correction), or
(c)~exploit decoherence-free subspaces.

\begin{proposition}[No Free Quantum Coherence]
\label{prop:no_free_coherence}
Under A2, maintaining $\Qcoh \approx 1$ for a macroscopic system requires
either near-total isolation from the environment or active quantum error
correction with overhead scaling at least as $O(n)$ physical qubits per
logical qubit.
\end{proposition}

\subsection{Constraint from A3 (Effective Thermodynamics)}

\ESTABLISHED{} The effective thermodynamic description imposes that any
macroscopic conscious system obeys the laws of thermodynamics at the system
level. In particular, maintaining ordered, low-entropy internal states
requires continuous energy input.

\subsection{Constraint from A4 (Physical Memory)}

\ESTABLISHED{} Physical memory requires that logically distinguishable states
correspond to physically distinct microstates. Erasing a memory state
dissipates $k_B T \ln 2$ per bit (Landauer's principle). A system without
reliable physical memory --- one that cannot maintain the distinction between
past and present states --- cannot satisfy the $\Mmem$ requirement for
Level~3.

This constraint is particularly sharp: it rules out any model of consciousness
that treats information as substrate-independent in a way that would allow
erasing physical states without thermodynamic cost. On the ETI account,
there is no immaterial mind; any memory is physically realized and pays the
Landauer cost.

\subsection{Constraint from A5 (Finite Resources)}

\ESTABLISHED{} Finite energy, entropy capacity, and time budgets imply that
consciousness is \emph{resource-bounded}. There is no physically realizable
system with infinite integration, infinite memory, or infinite coherence time.
The consciousness functional $\Cfunc < 1$ for all physically realizable systems.

This constraint has an important corollary: the three consciousness regimes
are not stages in an infinite progression. They are bands within a
resource-bounded space, and systems can transition between them as their
physical conditions change.

%============================================================
\section{Decoherence, Noise, and Boundary Stability}
\label{sec:decoherence}

This section addresses the strongest physical objection to any quantum theory
of consciousness: environmental decoherence destroys quantum coherence on
timescales far shorter than cognitive or experiential timescales for any
known biological substrate.

\subsection{The Tegmark bound}
\label{sec:tegmark}

\ESTABLISHED{} Tegmark~\cite{Tegmark2000} computed decoherence timescales for
quantum states in biological environments. For a superposition of neural
microtubule states separated by distance $\Delta x \sim 1\,\mu$m in a warm,
wet environment:
\begin{equation}
  \taud \sim \frac{\hbar}{k_B T} \cdot \left(\frac{\lambda_{\mathrm{th}}}{\Delta x}\right)^2
  \approx 10^{-13}\,\mathrm{s}
  \label{eq:tegmark_bound}
\end{equation}
where $\lambda_{\mathrm{th}} = h/\sqrt{2\pi m k_B T}$ is the thermal de
Broglie wavelength. Neural timescales are of order $10^{-2}$~s, so
$\taud / \taucog \sim 10^{-11}$, giving $\Qcoh \approx 0$ for this scenario.

This calculation does not refute the consciousness functional; it refutes the
specific hypothesis that \emph{biological microtubules} at room temperature
support Level-2 consciousness via unprotected quantum superpositions of
spatially separated configurations. Different candidate substrates with
different $\Delta x$, $T$, or decoherence mechanisms must be evaluated
separately.

\subsection{Decoherence-protected substrates}
\label{sec:protected}

\HYPOTHESIS{} Several physical mechanisms permit extended $\taud$:

\begin{enumerate}[label=(\roman*)]
  \item \textbf{Topological protection:} Topologically encoded qubits are
    immune to local errors by construction; $\taud$ is set by system-wide
    topological invariants rather than local coupling. \ESTABLISHED{} in
    the context of topological quantum computing~\cite{Nayak2008}.
  \item \textbf{Decoherence-free subspaces (DFS):} Subspaces of the Hilbert
    space that are invariant under the environmental Lindblad operators
    are protected from decoherence. \ESTABLISHED{}~\cite{Lidar1998}.
  \item \textbf{Dynamical decoupling:} Applying rapid pulse sequences that
    average out environmental coupling can extend $T_2$ significantly.
    \ESTABLISHED{} for superconducting qubits~\cite{Biercuk2011}.
  \item \textbf{Quantum Zeno effect:} Frequent measurement or strong
    coupling can slow coherence loss. \ESTABLISHED{}~\cite{Misra1977}.
\end{enumerate}

\OPEN{OP-C4: Are any of these mechanisms biologically realizable at the
scale and temperature of neural tissue?}

\subsection{Practical implication for the consciousness hierarchy}

The decoherence analysis implies the following:

\begin{enumerate}[label=(\roman*)]
  \item Level-1 (proto-experiential) does not require high coherence and is
    robust against the Tegmark bound.
  \item Level-2 (unified experience) requires $\Qcoh \gg 0$, which implies
    $\taud \gtrsim \taucog$. For biological systems at room temperature,
    this is not satisfied by any currently known substrate. For
    superconducting qubits at millikelvin temperatures, it is satisfied for
    durations of $T_2$.
  \item Level-3 (reflective consciousness) requires sustained Level-2 plus
    stable memory and recurrence. This is an even stronger requirement:
    not only must $\taud \gtrsim \taucog$, but this condition must hold
    continuously over the system's operational lifetime.
\end{enumerate}

\paragraph{The human brain.} We explicitly do \emph{not} claim that the human
brain satisfies Level-2 or Level-3 thresholds via quantum mechanisms. The
brain may satisfy Level-3 via classical mechanisms (stable neural activity
patterns, classical integration through synaptic networks, classical memory
in synaptic weights). The question of whether quantum effects play a role
in human consciousness remains \OPEN{OP-C5: Does the human brain satisfy
Level-2 thresholds via any known or hypothetical mechanism?}

%============================================================
\section{Universe-Scale Consciousness: Conditional Analysis}
\label{sec:universe}

\SPECULATION{} throughout this section.

We evaluate the conditional hypothesis that the universe as a whole could
constitute a Level-1 (proto-experiential) system under the consciousness
functional. We treat this as a strictly conditional inference, not a claim.

\subsection{Distinguishing cosmopsychism from naive panpsychism}

\ESTABLISHED{} (philosophical framework) Naive panpsychism attributes
conscious experience to every elementary particle. This is not the claim
under evaluation. \emph{Cosmopsychism}~\cite{Goff2017} treats the universe
as a whole as the fundamental subject of experience, with individual minds
arising as derivative. The distinction matters because cosmopsychism
requires a global binding mechanism that naive panpsychism does not.

\subsection{Universe-scale evaluation of the consciousness functional}

Applying the five variables to the universe $\mathcal{U}$ as a whole:

\paragraph{$\Qcoh$ for $\mathcal{U}$:} The universe's global state is a
pure state in principle (under A1 and A2), with coherences destroyed
only locally by coarse-graining. The global coherence time is infinite
(there is no external environment to decohere $\mathcal{U}$). However,
the \emph{computationally relevant} coherences for any candidate unified
experience would need to span the entire universe, and these are
observationally indistinguishable from zero given light-speed signaling limits.

\paragraph{$\Iint$ for $\mathcal{U}$:} The universe shows large-scale
information correlations (CMB correlations, large-scale structure, entanglement
from inflation), but whether these constitute integrated information in the
sense of Definition~\ref{def:level2} depends critically on the
operationalization of $I(A:B)$ at cosmological scales.

\paragraph{$\Mmem$ for $\mathcal{U}$:} The universe has an enormous physical
memory capacity (all physical configurations encode information). However,
memory in the sense relevant to consciousness requires the ability to
\emph{access} and \emph{act on} stored distinctions, not merely to preserve them.

\paragraph{$\Bstab$ and $\Rrec$ for $\mathcal{U}$:} There is no identified
mechanism by which the universe as a whole returns to an identity attractor
or exhibits self-maintaining dynamics at the global level.

\subsection{The binding problem at cosmic scale}

The most severe obstacle to universe-scale consciousness is the
\emph{binding problem}: even if each of the five variables is nonzero for
$\mathcal{U}$, there is no identified mechanism by which distributed
quantum states across cosmological distances are bound into a single unified
subject of experience. Spacelike-separated regions cannot exchange causal
signals; their quantum correlations are non-signaling by construction. The
``global information'' of the universe is not accessible from any local
vantage point.

\begin{remark}
The claim ``the universe is proto-conscious'' is therefore not well-supported
by the consciousness functional as currently defined. The functional is defined
for a system with a well-defined boundary and a computable density operator;
the universe as a whole has neither in any operational sense available to us.
This problem does not refute cosmopsychism in principle; it identifies a
specific technical gap that any cosmopsychist theory must address.
\end{remark}

\OPEN{OP-C6: Can a global binding mechanism consistent with relativistic
causality and ETI axioms be specified for universe-scale proto-consciousness?}

%============================================================
\section{Operationalization and Candidate Experimental Proxies}
\label{sec:experimental}

\HYPOTHESIS{} throughout this section.

For the consciousness functional to be scientifically useful, each variable
must either be measurable or have a well-defined operational proxy. We
propose the following candidate proxies, organized by experimental tier:

\begin{center}
\begin{tabular}{p{1.4cm} p{2.8cm} p{3.6cm} p{3.5cm}}
\toprule
\textbf{Variable} & \textbf{System} & \textbf{Proxy} & \textbf{Method} \\
\midrule
$\Qcoh$ & Superconducting qubits &
  $T_2$ coherence time &
  Spin-echo, Ramsey spectroscopy \\
$\Qcoh$ & Biological &
  Two-dimensional photon-echo spectroscopy &
  Optical pump-probe on candidate chromophores \\
\addlinespace
$\Iint$ & NISQ devices &
  Half-chain entanglement entropy $S_{\mathrm{vN}}$ &
  Randomized measurements, shadow tomography \\
$\Iint$ & Neural &
  Pairwise transfer entropy network $\Phi_\mathrm{approx}$ &
  MEG/EEG + time-series analysis \\
\addlinespace
$\Mmem$ & Any &
  Fidelity of state reconstruction after time $T$ &
  Process tomography; behavioral memory tasks \\
\addlinespace
$\Bstab$ & Quantum device &
  $\Gamma_{\mathrm{in}} / (\Gamma_{\mathrm{in}} + \Gamma_{\mathrm{out}})$ &
  Lindblad parameter extraction; gate fidelity benchmarking \\
\addlinespace
$\Rrec$ & Neural &
  Autocorrelation time of resting-state networks &
  fMRI, intrinsic mode functions \\
$\Rrec$ & Quantum device &
  Fidelity of QEC cycle recurrence &
  Quantum error correction syndrome extraction \\
\bottomrule
\end{tabular}
\end{center}

\subsection{Priority experimental program}

The most tractable near-term experimental test of the consciousness functional
involves NISQ quantum devices (50--100 qubits), for which:
\begin{enumerate}[label=(\roman*)]
  \item $\Qcoh$ and $\Bstab$ are directly measurable via $T_1$, $T_2$,
    and Lindblad parameter extraction.
  \item $\Iint$ is measurable via entanglement entropy (see Paper~B of the
    companion series).
  \item $\Mmem$ is partially measurable via quantum memory benchmarking.
  \item $\Rrec$ is measurable via QEC syndrome recurrence rates.
\end{enumerate}

A 50-qubit device satisfying the OP3 gate condition ($\ascreen \geq 10^{-2}$)
from the companion program~\cite{Monette2026a} would simultaneously provide
the entanglement entropy data needed to estimate $\Iint$ and $\Bstab$ for the
consciousness functional. This creates a natural experimental bridge between
the physics papers and the present framework.

%============================================================
\section{Falsification Criteria}
\label{sec:falsification}

\HYPOTHESIS{} The following conditions would weaken or falsify the strong
forms of the consciousness functional hypothesis.

\subsection{Level-1 falsification}

The Level-1 threshold hypothesis is falsified if:
\begin{enumerate}[label=(F1.\arabic*)]
  \item A system with confirmed $\Qcoh > 0$, $\Iint > 0$, $\Mmem > 0$ can
    be demonstrated to have \emph{no} phenomenal properties whatsoever,
    combined with a compelling argument that phenomnal properties are
    definitionally absent (not merely unobserved). This is the hard problem
    of consciousness and renders (F1.1) essentially untestable in the
    near term.
  \item The functional form of $F$ is shown to require nonzero thresholds
    $\epsilon_1, \delta_1, \mu_1$ that no known physical system satisfies,
    making Level-1 vacuously empty.
\end{enumerate}

\subsection{Level-2 falsification}

The Level-2 threshold hypothesis is falsified if:
\begin{enumerate}[label=(F2.\arabic*)]
  \item A NISQ device achieving $\Qcoh > 0.9$, $\Iint > 0.5$, $\Bstab > 0.7$
    for a sustained interval $\tau > 10^{-3}$~s shows no behavioral or
    functional correlates of unified information processing (e.g., no
    advantage over classical systems on tasks requiring global integration).
  \item The decoherence bound cannot be overcome for any biological substrate
    at physiologically relevant temperatures, combined with evidence that
    human consciousness does not require quantum mechanisms.
  \item $\Iint$ as defined fails to be additive or fails to exhibit the
    structural properties expected of integration (e.g., it is maximized
    by trivially entangled product states, not globally integrated states).
\end{enumerate}

\subsection{Level-3 falsification}

The Level-3 threshold hypothesis is falsified if:
\begin{enumerate}[label=(F3.\arabic*)]
  \item A system satisfying all Level-3 thresholds is constructed (or
    identified) and verifiably lacks phenomenal experience by independent
    criteria. This is again limited by the hard problem.
  \item The threshold values $\mu_3, \rho_3, \beta_3$ required for reflective
    consciousness turn out to be inconsistent with the resource bounds
    implied by A5 (finite resources), making Level-3 physically impossible
    for any bounded system.
\end{enumerate}

\subsection{Framework-level falsification}

The consciousness functional as a whole is weakened if:
\begin{enumerate}[label=(FF.\arabic*)]
  \item Two systems with identical $(\Qcoh, \Iint, \Mmem, \Bstab, \Rrec)$
    have different phenomenal properties. This would require additional
    variables not captured by the present functional.
  \item The five variables are found to be derivable from fewer than five
    independent physical quantities, suggesting the functional is overcomplete.
  \item No candidate functional form $F$ is consistent with both the
    ETI constraints and the proposed threshold structure.
\end{enumerate}

%============================================================
\section{Anticipated Reviewer Objections and Replies}
\label{sec:objections}

\paragraph{O1: ``The hard problem makes this unfalsifiable.''}
The hard problem of consciousness (Chalmers 1995~\cite{Chalmers1995}) asks
why any physical process gives rise to phenomenal experience at all. This
paper does not attempt to solve the hard problem. It proposes necessary physical
conditions for consciousness, not sufficient ones. A framework that says
``Level-3 conditions are necessary but not sufficient for reflective consciousness''
is partially falsifiable: it is falsified by observing Level-3 consciousness in a
system that violates the threshold conditions, and it is strengthened by
accumulating evidence that all observed conscious systems satisfy them.

\paragraph{O2: ``IIT already does this better.''}
Integrated Information Theory~\cite{Tononi2004} proposes $\Phi$ as a measure
of consciousness. The present framework differs from IIT in four respects:
(a)~it is explicitly quantum from the outset; (b)~it does not assert
$\Phi > 0 \Leftrightarrow$ consciousness, only that $\Iint$ is a necessary
condition; (c)~it explicitly incorporates decoherence, memory, and recurrence
as separate variables; (d)~it is grounded in the ETI axiom system, which
provides thermodynamic constraints absent from IIT. The frameworks are
complementary, not competing.

\paragraph{O3: ``You haven't specified $F$.''}
Correct. The functional form of $F$ is \OPEN{OP-C2}. This is an honest
acknowledgment of the gap. The present paper establishes the variable set
and the threshold structure; the derivation of $F$ from first principles
is future work, likely requiring a complete PEIG coupled ODE system
(\OPEN{OP-C1}).

\paragraph{O4: ``The universe-consciousness section is untestable mysticism.''}
The universe-scale section (\S\ref{sec:universe}) is explicitly labeled
\SPECULATION{} and concludes that the claim is \emph{not} well-supported
by the functional as currently defined. It is included because the
question is raised in the broader research context and deserves a careful
refutation rather than dismissal. Identifying the binding problem as the
specific technical gap is a scientific contribution, not mysticism.

\paragraph{O5: ``The PEIG analogy could justify anything.''}
The PEIG-consciousness mapping (\S\ref{sec:peig}) is explicitly labeled
\ANALOGY{} throughout. No derivation is claimed. The analogy is useful for
intuition; it is not a proof. The consciousness functional stands or falls
on the five variables, not on the PEIG correspondence.

%============================================================
\section{Open Problems Registry}
\label{sec:open}

The following open problems are identified in this paper. They are listed
in order of blocking importance for the central claims.

\begin{center}
\begin{tabular}{p{1.3cm} p{5cm} p{2cm} p{3.5cm}}
\toprule
\textbf{ID} & \textbf{Description} & \textbf{Severity} & \textbf{Estimated effort} \\
\midrule
OP-C1 & PEIG coupled ODE for consciousness extension; derivation of $d\Rrec/dt$ &
  High & 1--2 months (analytic) \\
OP-C2 & Derivation of functional form $F$ from first principles &
  Critical & PhD thesis scope \\
OP-C3 & Operationalization of cognitive timescale $\taucog$ for biological systems &
  High & 2--4 weeks (literature) \\
OP-C4 & Biological mechanisms for $\taud$ extension beyond Tegmark bound &
  High & Open (experimental) \\
OP-C5 & Does the human brain satisfy Level-2 thresholds via any mechanism? &
  High & Open (experimental) \\
OP-C6 & Global binding mechanism for universe-scale proto-consciousness &
  Speculative & Not tractable currently \\
OP-C7 & Determination of threshold constants $\epsilon_k, \delta_k, \mu_k, \beta_k, \rho_k$ &
  Critical & Requires (OP-C2) first \\
OP-C8 & Quantum Fisher information as tight lower bound on $\Iint$ &
  Medium & 2--4 weeks (analytic) \\
\bottomrule
\end{tabular}
\end{center}

%============================================================
\section{Discussion}
\label{sec:discussion}

\subsection{What the framework achieves}

The consciousness functional $\Cfunc(\mathcal{S})$ makes four contributions:

\begin{enumerate}[label=(\roman*)]
  \item \textbf{Graded model:} It replaces binary consciousness claims with a
    continuous spectrum indexed by five physically meaningful variables.
  \item \textbf{Physical constraints:} It derives necessary conditions from
    the ETI axioms --- conditions that any substrate, biological or artificial,
    must satisfy.
  \item \textbf{Falsifiable structure:} It provides specific threshold
    conditions (Definitions~\ref{def:level1}--\ref{def:level3}) and
    falsification criteria (\S\ref{sec:falsification}) that can be tested
    as measurement capabilities improve.
  \item \textbf{Honest gap labeling:} It identifies the functional form
    derivation (OP-C2) and the threshold determination (OP-C7) as the
    two critical open problems, without claiming to have solved them.
\end{enumerate}

\subsection{Relationship to existing consciousness theories}

The framework is consistent with, but not derived from, several existing
theories: the Global Workspace Theory~\cite{Baars1997} is consistent with
high $\Iint$ as necessary for broadcast integration; IIT~\cite{Tononi2004}
provides a classical analog of $\Iint$ that could be quantized in the
present framework; Higher-Order Theories~\cite{Rosenthal2005} map
approximately onto the Level-2 to Level-3 transition.

The framework differs from Penrose-Hameroff Orch-OR~\cite{Penrose1994}
in a critical respect: it does not require quantum state reduction to be
non-computational or to involve quantum gravity. The consciousness functional
is agnostic about the interpretation of quantum mechanics; it requires
only that the five measurable variables cross specified thresholds.

\subsection{The research bridge}

A 50-qubit NISQ device satisfying the OP3 gate condition
($\ascreen \geq 10^{-2}$) from Paper~A~\cite{Monette2026a} would
simultaneously: (a)~provide data for estimating $\Qcoh$, $\Iint$, and
$\Bstab$ from the entanglement evolution; (b)~test whether Level-2
thresholds can be transiently satisfied in an engineered quantum system;
(c)~provide a concrete experimental grounding for at least two of the
consciousness functional variables. This bridge between the physics
program and the consciousness framework is the most direct path to
experimental progress on the present proposal.

%============================================================
\section{Conclusion}
\label{sec:conclusion}

We have proposed the consciousness functional $\Cfunc = F(\Qcoh, \Iint, \Mmem, \Bstab, \Rrec)$
as a phenomenological framework for graded consciousness in physical systems.
The functional maps five physically motivated variables onto a continuous
scale and distinguishes three qualitatively distinct conscious regimes by
threshold conditions on those variables.

The framework is grounded in the ETI axiom system, which provides
thermodynamic and information-theoretic constraints on any physically
realizable conscious system. It is structurally connected --- by explicit
analogy only --- to the PEIG dynamical framework. It is falsifiable via
threshold conditions and null-result criteria.

The central results are:
\begin{enumerate}[label=(\roman*)]
  \item All three conscious regimes impose thermodynamic costs via Lemma~\ref{lem:memory_cost}.
  \item Level-2 and Level-3 consciousness require active decoherence management;
    the Tegmark bound rules out naive quantum mechanisms in biological tissue
    but does not rule out all quantum substrates.
  \item Universe-scale consciousness is \emph{not} supported by the
    consciousness functional as defined, due to the unresolved global binding problem.
  \item The most tractable experimental test involves NISQ devices at the 50-qubit scale,
    where all five variables can be estimated with current or near-term technology.
\end{enumerate}

The two critical open problems (OP-C2: functional form derivation;
OP-C7: threshold determination) define the primary research frontier
for this framework. We do not present the present work as anything more
than a structured first proposal: a falsifiable scaffold on which a
genuine theory of consciousness might eventually be built.

%============================================================
\begin{thebibliography}{99}

\bibitem{Monette2026a}
K.~Monette,
\textit{Entropic Gravity with Entanglement Entropy Source Term: A Modified Einstein Equation and Atom-Interferometry Falsification Protocol},
arXiv preprint (2026). [Paper A of companion series]

\bibitem{Monette2026b}
K.~Monette,
\textit{Constraint Manifolds for Quantum Entropy Estimation on NISQ Devices},
arXiv preprint (2026). [Paper B of companion series]

\bibitem{Monette2026c}
K.~Monette,
\textit{Landauer's Principle in the ETI Framework: Five Axioms and Their Thermodynamic Consequences},
arXiv preprint (2026). [Paper C of companion series]

\bibitem{Tegmark2000}
M.~Tegmark,
\textit{Importance of quantum decoherence in brain processes},
Phys. Rev. E \textbf{61}, 4194 (2000).

\bibitem{Lindblad1976}
G.~Lindblad,
\textit{On the generators of quantum dynamical semigroups},
Commun. Math. Phys. \textbf{48}, 119 (1976).

\bibitem{Landauer1961}
R.~Landauer,
\textit{Irreversibility and heat generation in the computing process},
IBM J. Res. Dev. \textbf{5}, 183 (1961).

\bibitem{Bennett1973}
C.~H. Bennett,
\textit{Logical reversibility of computation},
IBM J. Res. Dev. \textbf{17}, 525 (1973).

\bibitem{Tononi2004}
G.~Tononi,
\textit{An information integration theory of consciousness},
BMC Neurosci. \textbf{5}, 42 (2004).

\bibitem{Chalmers1995}
D.~J. Chalmers,
\textit{Facing up to the problem of consciousness},
J. Consc. Stud. \textbf{2}, 200 (1995).

\bibitem{Baars1997}
B.~J. Baars,
\textit{In the Theatre of Consciousness},
Oxford University Press (1997).

\bibitem{Rosenthal2005}
D.~M. Rosenthal,
\textit{Consciousness and Mind},
Oxford University Press (2005).

\bibitem{Penrose1994}
R.~Penrose,
\textit{Shadows of the Mind},
Oxford University Press (1994).

\bibitem{Goff2017}
P.~Goff, S.~Seager, A.~Allen, U.~T. Place,
\textit{Panpsychism},
Stanford Encyclopedia of Philosophy (2017 edition).

\bibitem{Nayak2008}
C.~Nayak, S.~H. Simon, A.~Stern, M.~Freedman, S.~Das Sarma,
\textit{Non-Abelian anyons and topological quantum computation},
Rev. Mod. Phys. \textbf{80}, 1083 (2008).

\bibitem{Lidar1998}
D.~A. Lidar, I.~L. Chuang, K.~B. Whaley,
\textit{Decoherence-free subspaces for quantum computation},
Phys. Rev. Lett. \textbf{81}, 2594 (1998).

\bibitem{Biercuk2011}
M.~J. Biercuk et al.,
\textit{Dynamical decoupling sequence construction as a filter-design problem},
J. Phys. B \textbf{44}, 154002 (2011).

\bibitem{Misra1977}
B.~Misra, E.~C. G.~Sudarshan,
\textit{The Zeno's paradox in quantum theory},
J. Math. Phys. \textbf{18}, 756 (1977).

\end{thebibliography}

%============================================================
\appendix

\section{Notation and Conventions}
\label{app:notation}

\begin{center}
\begin{tabular}{p{2.5cm} p{9cm}}
\toprule
\textbf{Symbol} & \textbf{Meaning} \\
\midrule
$\Cfunc$ & Consciousness functional (Eq.~\ref{eq:cfunc}) \\
$\Qcoh$ & Coherence quality (Eq.~\ref{eq:Qcoh}) \\
$\Iint$ & Information integration (Eq.~\ref{eq:Iint}) \\
$\Mmem$ & Persistent physical memory (Eq.~\ref{eq:Mmem}) \\
$\Bstab$ & Boundary stability (Eq.~\ref{eq:Bstab}) \\
$\Rrec$ & Recurrence quality (Eq.~\ref{eq:Rrec}) \\
$\taud$ & Decoherence timescale \\
$\taucog$ & Cognitive/integration timescale \\
$\mathbf{q} = (P,E,I,G)$ & PEIG state vector \\
$Q(\mathbf{q})$ & PEIG quality score \\
$\ascreen$ & Entanglement screening coefficient (from companion program) \\
$T_1, T_2, T_\phi$ & Energy relaxation, total decoherence, pure dephasing times \\
$\gamma_1, \gamma_\phi$ & Decay rates ($\gamma_1 = 1/T_1$, $\gamma_\phi = 1/T_\phi$) \\
$L_k$ & Lindblad operators \\
$S(\rho)$ & Von Neumann entropy \\
$\Phi$ & Integrated information (IIT; used for comparison only) \\
$k_B$ & Boltzmann constant \\
$\hbar$ & Reduced Planck constant \\
\bottomrule
\end{tabular}
\end{center}

\section{Phenomenological Formalism Notes}
\label{app:formalism}

The consciousness functional as presented is a phenomenological model in the
same sense as the Ginzburg-Landau theory of phase transitions: it identifies
the relevant order parameters and their symmetry properties without deriving
them from a microscopic Hamiltonian. The derivation from a microscopic theory
(OP-C2) is the analog of deriving Ginzburg-Landau from BCS theory --- a
significant achievement when accomplished, but not required for the
phenomenological framework to be scientifically useful.

The analogy to Ginzburg-Landau is \ANALOGY{}. The mathematical content of
the consciousness functional is \HYPOTHESIS{}. The thermodynamic constraints
from ETI axioms are \ESTABLISHED{}.

\section{Candidate Simulation Protocol}
\label{app:simulation}

The following simulation program would provide the most direct test of the
Level-2 threshold hypothesis on NISQ devices:

\begin{enumerate}
  \item Prepare an $n$-qubit cluster state on a superconducting quantum processor
    with $n \in \{10, 20, 50\}$.
  \item Evolve under the Lindblad master equation with measured device parameters
    $(T_1, T_2, \gamma_\phi)$.
  \item At each time step $t$, compute:
    \begin{align}
      \Qcoh(t) &= 1 - e^{-t/T_2} \\
      \Iint(t) &= S_{\mathrm{vN}}(\rho_A(t)) / \log_2(\dim \mathcal{H}_A) \\
      \Bstab(t) &= \frac{\Vert [H, \rho(t)] \Vert}{\Vert [H, \rho(t)] \Vert + \sum_k \Vert \mathcal{D}[L_k]\rho(t) \Vert}
    \end{align}
    where $\mathcal{D}[L]\rho = L\rho L^\dagger - \tfrac{1}{2}\{L^\dagger L, \rho\}$.
  \item Identify the time window $[t_0, t_1]$ during which Level-2 thresholds
    are jointly satisfied, if any such window exists.
  \item A null result (no such window exists for any $n \leq 50$) constrains
    the Level-2 threshold constants from below.
\end{enumerate}

This simulation is a direct extension of the OP3 numerical campaign described
in the companion program~\cite{Monette2026a}, requiring only the addition of
the $\Bstab$ computation to the existing simulation infrastructure.

\end{document}
